\documentclass[aspectratio=169,notheorems]{corvinusmetropolis}

\usepackage[hungarian]{babel}
\addbibresource{bibliography.bib}
\setbeamertemplate{frametitle continuation}{}

\renewcommand{\myinstlogo}{./corvinus-logos/color_simple_HUN.pdf}

\title{Biopénzügyek}
\subtitle{Biodiverzitás és Nature kockázat az európai piacon}
\author{Molnár Marcell\\Konzulens: Dr. Dömötör Barbara}
\email{marcell.molnar@stud.uni-corvinus.hu}

\institute[Corvinus University]
  {
  Pénzügy és számvitel alapszak\\
  Pénzügy specializáció\\
  Email: \putemail\\
  }
\date{2026. május 11.}
\titlegraphic{\hfill{\includegraphics[width=.2\textwidth]{\myinstlogo}}}
\keywords{biodiversity_finance, nature_finance, sustainable_finance}

\begin{document}
\simpleopen{}

\begin{frame}
  \frametitle{Irodalomi háttér}
  Az elmúlt években a Fenntartható Pénzügyek kutatói fokozott figyelmet fordított az új publikációkban a Biodiversity Finance ágnak\\ 
  Friss, releváns kutatások:
  \begin{itemize}
    \item \alert{biodiverzitási prémium} jelenléte és mértéke \parencite{coqueret_biodiversity_2024, naffa_biodiversity_2024}
    \item \alert{ENCORE adatbázis} biodiverzitás témakörben \parencite{clement_mapping_2025}
    \item \alert{private capital befektetések} \parencite{flammer_biodiversity_2025}\\
  \end{itemize}
  A hozzájárulásunk a szakirodalomhoz a biodiversity/nature hatások empirikus vizsgálata az európai részvénypiacon
\end{frame}

\begin{frame}
  \frametitle{Kutatási kérdések}
  \begin{enumerate}
    \item A biodiverzitási preferenciával rendelkező befektetők is az egész univerzumból választanak(long vagy short), nem pedig egyszerű kizárási stratégiát alkalmaznak.\footcite{pedersen_responsible_2021} 
    \item A piaci hozamokban megjelenik-e a biodiverzitási prémium\footcite{coqueret_biodiversity_2024} (ami elkülönül a carbon prémiumtól), és ha igen, akkor mekkora ennek a mértéke az európai piacon?\footcite{naffa_biodiversity_2024}
  \end{enumerate}
\end{frame}

\begin{frame}
  \frametitle{Adatbázisok, univerzum}
  Célunk az európai piac bemutatása, ehhez vagy a \alert{EURO STOXX 50} vagy az \alert{EURO STOXX 600} indexet szeretnénk használjuk.\\
   Egy másik lehetőség \alert{saját screening alapján} kiválasztani részvényeket az európai piacról, amik \alert{teljesítenek egy bizonyos feltételt}, például minimum piaci kapitalizáció.\\
  \medskip
  \metroset{block=fill}
  \begin{block}{Adatázisok}
    \begin{itemize}
    \item \alert{LSEG Refinitiv platform}: részvényárfolyamok, biodiverzitási mutatók, sustainability dashboard
    \item \alert{Bloomberg}: biodiverzitási és nature mutatók
      \item \alert{ENCORE adatbázis}: EU-ra iparágankénti nature dependency és impact mutatók
    \end{itemize}
  \end{block}
\end{frame}

\bibframes{Hivatkozások}
\end{document}
